\documentclass[a4paper,11pt,twoside]{StyleThese}

\usepackage{amsmath,amssymb}             % AMS Math
\usepackage[french]{babel}
\usepackage[latin1]{inputenc}
\usepackage[T1]{fontenc}
\usepackage[left=1.5in,right=1.3in,top=1.1in,bottom=1.1in,includefoot,includehead,headheight=13.6pt]{geometry}
\renewcommand{\baselinestretch}{1.05}

% Table of contents for each chapter

\usepackage[nottoc, notlof, notlot]{tocbibind}
\usepackage[french]{minitoc}
\setcounter{minitocdepth}{2}
\mtcindent=15pt
% Use \minitoc where to put a table of contents

\usepackage{aecompl}

% Glossary / list of abbreviations

\usepackage[intoc]{nomencl}
\renewcommand{\nomname}{Liste des Abr�viations}

\makenomenclature

% My pdf code

\usepackage{ifpdf}

\ifpdf
  \usepackage[pdftex]{graphicx}
  \DeclareGraphicsExtensions{.jpg}
  \usepackage[a4paper,pagebackref,hyperindex=true]{hyperref}
\else
  \usepackage{graphicx}
  \DeclareGraphicsExtensions{.ps,.eps}
  \usepackage[a4paper,dvipdfm,pagebackref,hyperindex=true]{hyperref}
\fi

\graphicspath{{.}{images/}}

%nicer backref links
\renewcommand*{\backref}[1]{}
\renewcommand*{\backrefalt}[4]{%
\ifcase #1 %
(Non cit�.)%
\or
(Cit� en page~#2.)%
\else
(Cit� en pages~#2.)%
\fi}
\renewcommand*{\backrefsep}{, }
\renewcommand*{\backreftwosep}{ et~}
\renewcommand*{\backreflastsep}{ et~}

% Links in pdf
\usepackage{color}
\definecolor{linkcol}{rgb}{0,0,0.4} 
\definecolor{citecol}{rgb}{0.5,0,0} 

% Change this to change the informations included in the pdf file

\hypersetup
{
bookmarksopen=true,
pdftitle="Cr�ation et utilisation d'atlas anatomiques num�riques pour la radioth�rapie",
pdfauthor="Olivier COMMOWICK", %auteur du document
pdfsubject="Segmentation d'images par atlas et cr�ation d'atlas", %sujet du document
%pdftoolbar=false, %barre d'outils non visible
pdfmenubar=true, %barre de menu visible
pdfhighlight=/O, %effet d'un clic sur un lien hypertexte
colorlinks=true, %couleurs sur les liens hypertextes
pdfpagemode=None, %aucun mode de page
pdfpagelayout=SinglePage, %ouverture en simple page
pdffitwindow=true, %pages ouvertes entierement dans toute la fenetre
linkcolor=linkcol, %couleur des liens hypertextes internes
citecolor=citecol, %couleur des liens pour les citations
urlcolor=linkcol %couleur des liens pour les url
}

% definitions.
% -------------------

\setcounter{secnumdepth}{3}
\setcounter{tocdepth}{2}

% Some useful commands and shortcut for maths:  partial derivative and stuff

\newcommand{\pd}[2]{\frac{\partial #1}{\partial #2}}
\def\abs{\operatorname{abs}}
\def\argmax{\operatornamewithlimits{arg\,max}}
\def\argmin{\operatornamewithlimits{arg\,min}}
\def\diag{\operatorname{Diag}}
\newcommand{\eqRef}[1]{(\ref{#1})}

\usepackage{rotating}                    % Sideways of figures & tables
%\usepackage{bibunits}
%\usepackage[sectionbib]{chapterbib}          % Cross-reference package (Natural BiB)
%\usepackage{natbib}                  % Put References at the end of each chapter
                                         % Do not put 'sectionbib' option here.
                                         % Sectionbib option in 'natbib' will do.
\usepackage{fancyhdr}                    % Fancy Header and Footer

% \usepackage{txfonts}                     % Public Times New Roman text & math font
  
%%% Fancy Header %%%%%%%%%%%%%%%%%%%%%%%%%%%%%%%%%%%%%%%%%%%%%%%%%%%%%%%%%%%%%%%%%%
% Fancy Header Style Options

\pagestyle{fancy}                       % Sets fancy header and footer
\fancyfoot{}                            % Delete current footer settings

%\renewcommand{\chaptermark}[1]{         % Lower Case Chapter marker style
%  \markboth{\chaptername\ \thechapter.\ #1}}{}} %

%\renewcommand{\sectionmark}[1]{         % Lower case Section marker style
%  \markright{\thesection.\ #1}}         %

\fancyhead[LE,RO]{\bfseries\thepage}    % Page number (boldface) in left on even
% pages and right on odd pages
\fancyhead[RE]{\bfseries\nouppercase{\leftmark}}      % Chapter in the right on even pages
\fancyhead[LO]{\bfseries\nouppercase{\rightmark}}     % Section in the left on odd pages

\let\headruleORIG\headrule
\renewcommand{\headrule}{\color{black} \headruleORIG}
\renewcommand{\headrulewidth}{1.0pt}
\usepackage{colortbl}
\arrayrulecolor{black}

\fancypagestyle{plain}{
  \fancyhead{}
  \fancyfoot{}
  \renewcommand{\headrulewidth}{0pt}
}

\usepackage{MyAlgorithm}
\usepackage[noend]{MyAlgorithmic}

%%% Clear Header %%%%%%%%%%%%%%%%%%%%%%%%%%%%%%%%%%%%%%%%%%%%%%%%%%%%%%%%%%%%%%%%%%
% Clear Header Style on the Last Empty Odd pages
\makeatletter

\def\cleardoublepage{\clearpage\if@twoside \ifodd\c@page\else%
  \hbox{}%
  \thispagestyle{empty}%              % Empty header styles
  \newpage%
  \if@twocolumn\hbox{}\newpage\fi\fi\fi}

\makeatother
 
%%%%%%%%%%%%%%%%%%%%%%%%%%%%%%%%%%%%%%%%%%%%%%%%%%%%%%%%%%%%%%%%%%%%%%%%%%%%%%% 
% Prints your review date and 'Draft Version' (From Josullvn, CS, CMU)
\newcommand{\reviewtimetoday}[2]{\special{!userdict begin
    /bop-hook{gsave 20 710 translate 45 rotate 0.8 setgray
      /Times-Roman findfont 12 scalefont setfont 0 0   moveto (#1) show
      0 -12 moveto (#2) show grestore}def end}}
% You can turn on or off this option.
% \reviewtimetoday{\today}{Draft Version}
%%%%%%%%%%%%%%%%%%%%%%%%%%%%%%%%%%%%%%%%%%%%%%%%%%%%%%%%%%%%%%%%%%%%%%%%%%%%%%% 

\newenvironment{maxime}[1]
{
\vspace*{0cm}
\hfill
\begin{minipage}{0.5\textwidth}%
%\rule[0.5ex]{\textwidth}{0.1mm}\\%
\hrulefill $\:$ {\bf #1}\\
%\vspace*{-0.25cm}
\it 
}%
{%

\hrulefill
\vspace*{0.5cm}%
\end{minipage}
}

\let\minitocORIG\minitoc
\renewcommand{\minitoc}{\minitocORIG \vspace{1.5em}}

\usepackage{multirow}
\usepackage{slashbox}

\newenvironment{bulletList}%
{ \begin{list}%
	{$\bullet$}%
	{\setlength{\labelwidth}{25pt}%
	 \setlength{\leftmargin}{30pt}%
	 \setlength{\itemsep}{\parsep}}}%
{ \end{list} }

\newtheorem{definition}{D�finition}
\renewcommand{\epsilon}{\varepsilon}

% centered page environment

\newenvironment{vcenterpage}
{\newpage\vspace*{\fill}\thispagestyle{empty}\renewcommand{\headrulewidth}{0pt}}
{\vspace*{\fill}}



\begin{document}

\begin{titlepage}
\begin{center}
\noindent {\large \textbf{UNIVERSIT� DE NICE - SOPHIA ANTIPOLIS}} \\
\vspace*{0.3cm}
\noindent {\LARGE \textbf{�COLE DOCTORALE STIC}} \\
\noindent \textbf{SCIENCES ET TECHNOLOGIES DE L'INFORMATION \\ ET DE LA COMMUNICATION} \\
\vspace*{0.5cm}
\noindent \Huge \textbf{T H � S E} \\
\vspace*{0.3cm}
\noindent \large {pour obtenir le titre de} \\
\vspace*{0.3cm}
\noindent \LARGE \textbf{Docteur en Sciences} \\
\vspace*{0.3cm}
\noindent \Large de l'Universit� de Nice - Sophia Antipolis \\
\noindent \Large \textbf{Mention : \textsc{Informatique}}\\
\vspace*{0.4cm}
\noindent \large {Pr�sent�e et soutenue par\\}
\noindent \LARGE Olivier \textsc{Commowick} \\
\vspace*{0.8cm}
\noindent {\Huge \textbf{Cr�ation et utilisation d'atlas anatomiques num�riques pour la radioth�rapie}} \\
\vspace*{0.8cm}
\noindent \Large Th�se dirig�e par Gr�goire \textsc{Malandain} \\
\vspace*{0.2cm}
\noindent \Large pr�par�e � l'INRIA Sophia Antipolis, Projet \textsc{Asclepios} \\
\vspace*{0.2cm}
\noindent \large soutenue le 14 f�vrier 2007 \\
\vspace*{0.5cm}
\end{center}
\noindent \large \textbf{Jury :} \\
\begin{center}
\noindent \large 
\begin{tabular}{llcl}
      \textit{Rapporteurs :}	& Patrick \textsc{Clarysse}		& - & CNRS (CREATIS)\\
				& Louis \textsc{Collins}		& - & McGill University\\
      \textit{Directeur :}	& Gr�goire \textsc{Malandain}		& - & INRIA (Asclepios)\\
      \textit{Pr�sident :}	& Nicholas \textsc{Ayache}		& - & INRIA (Asclepios)\\
      \textit{Examinateurs :}   & Pierre-Yves \textsc{Bondiau}          & - & Centre Antoine Lacassagne (Nice)\\
      				& Guido \textsc{Gerig}			& - & University of North Carolina\\
      				& Vincent \textsc{Gr�goire}		& - & Universit� Catholique de Louvain\\
      \textit{Invit� :}		& Hanna \textsc{Kafrouni}		& - & DOSISoft S.A.
\end{tabular}
\end{center}
\end{titlepage}
\sloppy

\titlepage


\dominitoc

\pagenumbering{roman}

 \cleardoublepage

\section*{Remerciements}

A faire en dernier :-)

\tableofcontents

\mainmatter

\chapter{Introduction}
\label{chap:intro}
\minitoc

\section{Exemple d'illustration}

\subsection{Une sous section pour le fun}

Eh non je n'�crirai pas votre th�se... D�sol�. Juste pour montrer que ce style g�re les accents et comment on inclut une figure. A noter, les path vers les dossiers d'images et l'extension sont d�finis dans formatAndDefs.tex . Pour pdflatex, les images jpg sont g�r�es. Pour latex, utiliser ps ou eps.

\begin{figure}[!htbp]
  \begin{center}
    \includegraphics[width=0.9\textwidth]{Chapitre1/arctic_control}
  \end{center}
  \caption{Une jolie image...}
  \label{fig:jolieImage}
\end{figure}

\section{Une �quation}

Juste pour montrer les commandes argmin et d�riv�e partielle.

\begin{equation}
  T = \argmin_T E(T,R,F)
\end{equation}

R�gularisation:

\begin{equation}
  \pd{T}{t} = \Delta T
\end{equation}

\section{Une autre section}

Pour montrer l'environnement bulletList

\begin{bulletList}
 \item Premier point
 \item Deuxi�me point
% \item Et une r�f�rence ajout�e dans la liste des abr�viations \nomenclature{DTI}{Diffusion Tensor Imaging} DTI
\end{bulletList}


\appendix

\chapter{Exemple d'annexe}
\label{chap:annexe1}

\section{Exemple d'annexe}

Et j'en profite pour me citer et montrer mon style bibtex (deux auteurs) \cite{Commowick_MICCAI_2007}.

Et deux autres citations pour le style optionnel de bibliographie : un seul auteur \cite{Oakes_RStat_1999} et plus de deux auteurs \cite{Guimond_CVIU_2000}.

\bibliographystyle{StyleThese}
\bibliography{These}

%\printnomenclature

\cleardoublepage
\begin{vcenterpage}
\noindent\rule[2pt]{\textwidth}{0.5pt}
\\
{\large\textbf{R�sum� :}}
L'objectif de cette th�se est de fournir aux radioth�rapeutes des outils de contourage automatique des structures � risque pour la planification de la radioth�rapie des tumeurs c�r�brales et de la r�gion ORL.
\\
Nous utilisons pour cela un atlas anatomique, constitu� d'une repr�sentation de l'anatomie associ�e � une image de celle-ci. Le recalage de cet atlas permet de contourer automatiquement les organes du patient et ainsi obtenir un gain de temps consid�rable. Les contributions pr�sent�es se concentrent sur trois axes.
\\
Tout d'abord, nous souhaitons obtenir une m�thode de recalage la plus ind�pendante possible du r�glage de ses param�tres. Celui-ci, effectu� par le m�decin, se doit d'�tre minimal, tout en garantissant un r�sultat robuste. Nous proposons donc des m�thodes de recalage permettant un meilleur contr�le de la transformation obtenue, en passant par des techniques de rejet d'appariements aberrants ou en utilisant des transformations localement affines.
\\
Le second axe est consacr� � la prise en compte de structures dues � la tumeur. En effet, la pr�sence de ces structures, absentes de l'atlas, perturbe le recalage de celui-ci. Nous proposons donc �galement des m�thodes afin de contourer ces structures et de les prendre en compte dans le recalage.
\\
Enfin, nous pr�sentons la construction d'un atlas ORL et son �valuation sur une base de patients. Nous montrons ici la faisabilit� de l'utilisation d'un atlas de cette r�gion, ainsi qu'une m�thode simple afin d'�valuer les m�thodes de recalage utilis�es pour construire un atlas.
\\
L'ensemble de ces travaux a �t� impl�ment� dans le logiciel Imago de DOSIsoft, ceci ayant permis d'effectuer une validation en conditions cliniques.
\\
\\
{\large\textbf{Mots cl�s :}}
Segmentation par atlas, recalage non lin�aire, radioth�rapie, cr�ation d'atlas
\\
\noindent\rule[2pt]{\textwidth}{0.5pt}
\end{vcenterpage}

\begin{vcenterpage}
\noindent\rule[2pt]{\textwidth}{0.5pt}
\begin{center}
{\large\textbf{Design and Use of Numerical Anatomical Atlases for Radiotherapy\\}}
\end{center}
{\large\textbf{Abstract:}}
The main objective of this thesis is to provide radio-oncology specialists with automatic tools for delineating organs at risk of a patient undergoing a radiotherapy treatment of cerebral or head and neck tumors.
\\
To achieve this goal, we use an anatomical atlas, i.e. a representative anatomy associated to a clinical image representing it. The registration of this atlas allows to segment automatically the patient structures and to accelerate this process. Contributions in this method are presented on three axes.
\\
First, we want to obtain a registration method which is as independent as possible w.r.t. the setting of its parameters. This setting, done by the clinician, indeed needs to be minimal while guaranteeing a robust result. We therefore propose registration methods allowing to better control the obtained transformation, using outlier rejection techniques or locally affine transformations.
\\
The second axis is dedicated to the consideration of structures associated with the presence of the tumor. These structures, not present in the atlas, indeed lead to local errors in the atlas-based segmentation. We therefore propose methods to delineate these structures and take them into account in the registration.
\\
Finally, we present the construction of an anatomical atlas of the head and neck region and its evaluation on a database of patients. We show in this part the feasibility of the use of an atlas for this region, as well as a simple method to evaluate the registration methods used to build an atlas.
\\
All this research work has been implemented in a commercial software (Imago from DOSIsoft), allowing us to validate our results in clinical conditions.
\\
\\
{\large\textbf{Keywords:}}
Atlas-based Segmentation, non rigid registration, radiotherapy, atlas creation
\\
\noindent\rule[2pt]{\textwidth}{0.5pt}
\end{vcenterpage}

\end{document}
